% titel: Rotierende Scheiben: gemeinsame Winkelgeschwindigkeit
% skript: Kap. Rotation, Drehimpulserhaltung, Reibungsenergie

Zwei Scheiben rotieren zunächst getrennt um dieselbe (vertikale) Achse; durch
Aufsetzen der oberen auf die untere wirken zwischen ihnen Reibungskräfte, bis
beide mit gemeinsamer Winkelgeschwindigkeit drehen. Ohne äussere Drehmomente
um die Achse bleibt der \emph{Drehimpuls} erhalten (nicht die Energie).

\begin{werkzeugkasten}[Werkzeuge]
\textbf{Drehimpulserhaltung} für die Achse (keine äusseren Drehmomente):
\begin{equation}
  L_\mathrm{vor} = L_\mathrm{nach}
  \;\Rightarrow\;
  I_1\omega_1 + I_2\omega_2 = (I_1+I_2)\,\omega\;.
\end{equation}

\textbf{Trägheitsmoment} einer homogenen Scheibe der Masse $m$ und des Radius
$R$ um die Symmetrieachse:
\begin{equation}
  I = \frac12 m R^2\;.
\end{equation}

\textbf{Rotationsenergie:} $E_\mathrm{rot}=\frac12 I\omega^2$.

\textbf{Reibungsenergie:} Der Energieverlust beim Reibschluss ist die Differenz
der Rotationsenergien vor und nach dem Vorgang.
\end{werkzeugkasten}

\begin{werkzeugkasten}[Vorgehen in allgemeiner Form]
$I_1,I_2$ aus $\frac12 mR^2$; aus $I_1\omega_1+I_2\omega_2=(I_1+I_2)\omega$
die gemeinsame $\omega$; Reibungsenergie $=\frac12 I_1\omega_1^2
+\frac12 I_2\omega_2^2-\frac12 (I_1+I_2)\omega^2$.
\end{werkzeugkasten}

\begin{achtung}[Typische Fallen]
\begin{itemize}
  \item Beim Schleifen ist die Gesamtenergie nicht erhalten - nur der
       Drehimpuls. Die Reibung setzt Energie in Wärme um.
  \item $\omega_2<0$: Vorzeichen der zweiten Scheibe (Gegendrehung) sauber
       in die Erhaltungsgleichung nehmen.
  \item $60\,\mathrm{min}^{-1}$ sind $1\,\mathrm s^{-1}$ (Umdrehungen),
       also $\omega=2\pi\,\mathrm{rad/s}$; nicht verwechseln.
\end{itemize}
\end{achtung}